\section{歴史的背景と問題意識 - 古典物理学の限界}

\noindent\textbf{この章の中心命題}
\[
\boxed{
u_{\mathrm{RJ}}(\nu,T)\xrightarrow[\nu\to\infty]{}\infty,\qquad
K_{\max}=h\nu-\Phi,\qquad
h\nu=E_n-E_m
}
\]
黒体放射、光電効果、原子スペクトルの三つの事実は、微視的世界では古典力学の状態概念と連続的なエネルギー交換の仮定をそのまま使えないことを示す。

19世紀末までに、ニュートン力学とマクスウェル電磁気学によって、巨視的な物理現象の大部分は説明されたと考えられていた。したがって、もし微視的世界も同じ枠組みで記述できるなら、粒子の状態は古典的な位相空間上の点として表され、時間発展も古典運動方程式で決まるはずである。しかし実際には、微視的現象ではこの期待が破れる。
\begin{itemize}
\item{黒体放射と紫外線発散}: 
古典論では、振動数 $\nu$ から $\nu+d\nu$ の間にある電磁場モード数と等分配則を用いると、単位体積・単位振動数あたりのエネルギー密度は
\[
u_{\mathrm{RJ}}(\nu,T)=\frac{8\pi \nu^2}{c^3}k_{\mathrm{B}}T
\]
となる。したがって全エネルギー密度は
\[
\int_0^\infty u_{\mathrm{RJ}}(\nu,T)\,d\nu = \infty
\]
と発散してしまう。しかし実際に観測されるスペクトルは高振動数で減衰しており、
\[
u_{\mathrm{P}}(\nu,T)=\frac{8\pi h\nu^3}{c^3}\frac{1}{e^{h\nu/(k_{\mathrm{B}}T)}-1}
\]
でよく記述される。この不一致が、エネルギー量子 $E=h\nu$ の導入を要請した。

\item{光電効果}: 
古典的な波動論だけで考えれば、電子へ渡されるエネルギーは光の強度に比例し、十分長く照射すればどの振動数でも電子は飛び出せるはずである。しかし実験では、電子が飛び出すためにはしきい振動数 $\nu_0$ が必要であり、飛び出した電子の最大運動エネルギーは
\[
K_{\max}=h\nu-\Phi
\]
に従う。ここで $\Phi$ は仕事関数である。すなわち、電子放出の有無を決めるのは強度ではなく振動数であり、古典的波動像だけでは説明できない。

\item{原子の安定性の問題}: 
古典電磁気学では、電荷をもつ粒子が加速するとラーモアの式
\[
P=\frac{e^2 a^2}{6\pi\varepsilon_0 c^3}
\]
に従って電磁波を放射する。したがって原子核のまわりを回る電子はエネルギーを失い、やがて原子核へ落ち込むはずである。しかし実際には原子は安定であり、しかもスペクトルは連続ではなく
\[
h\nu = E_n-E_m
\]
の形の離散線として観測される。ここでも古典理論の予測は観測と一致しない。
\end{itemize}
以上の事情から、古典力学における状態の記述および時間発展の法則をそのまま微視的世界へ持ち込むことはできない。したがって必要なのは、観測事実と両立するように状態・観測量・時間発展を新しく定める理論である。
